% Einheit 4 von 4 – Lagebefunde im Sachzusammenhang (bank/lagebeziehungen/e4.jsonl, zusammenbau v0.1)
%% TODO Zeitmarke Einheit 4: Marken-Zeile nicht lesbar
%% TODO Zweigzeile Teil 1: was hier gelernt wird (Satz in Schülersprache aus dem Einheitstitel) – Einheit 4
\einheitenkopf[e4]{Einheit 4 von 4 · Lagebefunde im Sachzusammenhang}
\zweigzeile{keine Prüfungsaufgabe}
\setcounter{aufgabe}{15}

%% TODO Titel: Ich-kann-Satz fehlt in der Bank (Platzhalter: Kettenname) – Nr. 16
%% TODO Anweisung: ein Satz über der ersten Teilaufgabe fehlt in der Bank – Nr. 16
\begin{aufgabe}{Lagebefunde im Sachzusammenhang}
%% TODO \rechenplatz{n}: Schrittzahl des Grundfalls fehlt in der Bank (nur bei fester Schreibform, 2.3 b)
\begin{teile}
\teil Frage: Liegt der Schatten der Lampe auf der Wand? Rechnung: Durchstoßpunkt der Lichtgeraden mit der Wandebene $(2 | 0 | 1{,}8)$. Ist die Frage damit beantwortet? \\ \kreuz{ja} \\ \kreuz{nein – ein Bereich ist noch zu prüfen}
\teil Eine Wand liegt in der Ebene $y = 6$ mit $0 \le x \le 6$ und $0 \le z \le 3$. Eine Lampe in $L(2 | 0 | 4)$ beleuchtet die Spitze $P(3 | 2 | 3)$ eines Pfostens. Liegt der Schatten von $P$ auf der Wand?
\teil Sonnenlicht fällt in Richtung $(-1 | -2 | -1)$; eine Hauswand liegt in der xz-Ebene mit $0 \le x \le 6$ und $0 \le z \le 3$. Schritt I: $(x | 0 | z) = (4 | 2 | 3) + \lambda \cdot (-1 | -2 | -1)$ liefert $\lambda = 1$ und $(3 | 0 | 2)$. Schritt II: $0 < 3 < 6$ und $0 < 2 < 3$. Erläutere beide Schritte im Sachzusammenhang. \hfill (Abitur 2024 GK)
\teil Ein Netz steht in der Ebene $y = 6$, seine Oberkante ist $2$ m hoch (1 LE = 1 m, die xy-Ebene ist der Boden). Ein Ball fliegt von $P(1 | 4 | 3)$ geradlinig in Richtung $(0 | 4 | -1)$. Untersuche rechnerisch, ob er das Netz berührt. \hfill (Abitur 2018 LK)
\teil Ein Spielfeld ist das Rechteck mit $0 \le x \le 9$ und $0 \le y \le 18$ auf dem Boden (xy-Ebene, 1 LE = 1 m). Nach dem Schlag bewegt sich der Ball auf $X_t(4 | 8t + 3 | -5t^2 + 3t + 2)$, $t$ in Sekunden. Untersuche, ob der Ball im Spielfeld auf den Boden trifft. \hfill (Abitur 2018 LK)
\end{teile}
\end{aufgabe}

%% TODO Titel: Ich-kann-Satz fehlt in der Bank (Platzhalter: Kettenname) – Nr. 17
%% TODO Anweisung: ein Satz über der ersten Teilaufgabe fehlt in der Bank – Nr. 17
\begin{aufgabe}{Gerade und Ebene: Länge eines Schattens auf einer Dachfläche über Schnitt von Lichtgerade und Ebene beschreiben}
\begin{teile}
\teil Auf dem schrägen Dach eines Carports, das in der Ebene $D$ liegt, steht im Punkt $F$ eine senkrechte Antenne mit der Spitze $T$. Das Sonnenlicht fällt in parallelen Geraden mit dem Richtungsvektor $\vec{v}$; der Schatten der Antenne liegt ganz auf dem Dach. Beschreibe, wie man die Länge des Schattens berechnet, wenn $T$ und $\vec{v}$ bekannt sind. \hfill (Abitur 2017 LK)
\end{teile}
\end{aufgabe}

%% TODO Titel: Ich-kann-Satz fehlt in der Bank (Platzhalter: Kettenname) – Nr. 18
%% TODO Anweisung: ein Satz über der ersten Teilaufgabe fehlt in der Bank – Nr. 18
\begin{aufgabe}{Lagebefunde im Sachzusammenhang – Fehler finden, Begründen}
\begin{teile}
\teil Ich finde den Fehler: Eine Wand liegt in der Ebene $y = 4$ mit $0 \le x \le 1{,}8$ und $0 \le z \le 2{,}5$. Eine Lampe in $L(0 | 0 | 5)$ beleuchtet $P(1 | 2 | 4)$. Paula berechnet den Durchstoßpunkt $(2 | 4 | 3)$ und schreibt: Der Schatten liegt auf der Wand.
\teil Begründe, warum der Schatten eines Punkts $P$ auf einer Wand der Durchstoßpunkt der Lichtgeraden durch $P$ mit der Wandebene ist.
\end{teile}
\end{aufgabe}
